% Solutions to the Chapter 5 exercises that are not code, from lectures/L05/appendix/c_exercises.md.
% Exercises 5.6 and 5.7 are Code and are not answered; see chapter.tex for why.

\section{Chapter 5: Character Device Drivers}
\label{sol:sec:c5}
Answers to the exercises in \secref{c5:sec:exercises}. Exercises~\ref{c5:ex:6} and~\ref{c5:ex:7}
are Code, and are checked on the target as their Check yourself lines say.

% ------------------------------------------------------------------------------------------
\solution[fillaccent2]{c5:ex:1}{The node and the number}{Recall}

\xpart{a} The number is whatever was free when the driver loaded, so the script breaks in two
separate ways. \textbf{The driver moves:} a new module, a different load order or a kernel update
that registers another dynamically numbered driver first, and \code{qa_fifo} gets a different
major; the script then fails to find its device. \textbf{Something else takes 243:} after the same
kind of change another driver can be handed 243, and the script opens that driver's device and
talks to it without any error at all. The second is the worse of the two, because nothing fails.

\xpart{b} It creates a second name for the same device: a directory entry of type character with
major 243 and minor 0, written without consulting anything (\secref{c5:app:a1}). Opening it opens
exactly what opening \file{/dev/qa_fifo} opens, the same driver and the same FIFO, so bytes written
through one name come out of the other.

\xpart{c} Step 5, \code{device_create}. \file{/proc/devices} lists the name as soon as the number
is allocated (step 2), but the node is made by devtmpfs in response to the uevent that
\code{device_create} emits (\secref{c5:app:a6}); no device, no uevent, no node. If
\code{class_create} was skipped as well, there was no class to create the device in.

\xpart{d} \code{ENXIO} comes from the kernel's lookup of a registered \code{cdev} for the node's
major and minor, before any driver is involved: nothing is registered for that pair. Either the
driver is not loaded, or it allocated a number but never called \code{cdev_add}, or the node is
stale and carries numbers from an earlier boot. It rules out a missing node (\code{ENOENT}), a
permissions problem (\code{EACCES}), and anything in your own \code{open}, which was never called.

% ------------------------------------------------------------------------------------------
\solution[fillaccent2]{c5:ex:2}{The table}{Recall}

\xpart{a} \code{.owner} is what makes \code{open} take a reference on the module, and the reference
is what makes \code{rmmod} refuse with ``Module is in use''. Without it the module's use count stays
at zero while a process holds an open file whose \code{f_op} points at the module's
\code{file_operations} table. \code{rmmod} succeeds and frees the module's code and data, table
included. The process's next \code{read}, \code{write} or \code{close} calls through a pointer into
freed memory: an oops if you are lucky, and whatever has been loaded there since if you are not.

\xpart{b} Not a bug, as long as there is no per-open state to free, which is the case here (the
table in \secref{c5:app:b2} marks \code{.release} as not required). Nothing is called instead:
when the last reference to the file goes, the VFS calls \code{.release} only if it is set, and
otherwise simply drops its own reference on the module and frees the \code{struct file}.

\xpart{c} Yes. A missing \code{.open} is not a refusal: the character device layer calls the
driver's \code{open} only if there is one, and otherwise the open succeeds and returns a
descriptor.

\xpart{d} The device is live. Any process that opens a node with your numbers, and a node can exist
before \code{device_create} makes one, can call your \code{open}, \code{read} and \code{write}
while the rest of your init function is still running, and on another CPU at the same time. The
rule: \textbf{everything a callback touches must be fully set up before the call that publishes it},
so the FIFO is created before \code{cdev_add}, never after (\secref{c5:app:b2}).

\xpart{e} It is called when the last reference to the open file goes, which is when the last
descriptor referring to it is closed. The simplest case: a process opens the device and then calls
\code{fork}. Parent and child now hold descriptors for the same open file. The parent's
\code{close} drops one reference and calls nothing; the child's \code{close} drops the last and
calls \code{release}. \code{dup} followed by two \code{close}s in one process does the same.

% ------------------------------------------------------------------------------------------
\solution[fillaccent3]{c5:ex:3}{Return values}{Hand calculation}

\xpart{a} 40. \code{read} returns 40, the buffer holds the 40 bytes, and the FIFO is now empty. A
short count is a success.

\xpart{b} \code{-EAGAIN}. \code{read} returns $-1$ with \code{errno} set to \code{EAGAIN},
``Resource temporarily unavailable''.

\xpart{c} 0. A count of 0 returns 0 and consumes nothing, so the level stays 40. The program sees
the same 0 that means end of file everywhere else, which is harmless only because it asked for
nothing.

\xpart{d} 60. The first 60 bytes are accepted and the FIFO is full; the program sees 60 and must
write the other 40 later.

\xpart{e} \code{-ENOSPC}. \code{write} returns $-1$ with \code{errno} set to \code{ENOSPC}, ``No
space left on device''.

\xpart{f} \code{-EFAULT}, provided there are bytes to deliver: the driver takes them from the FIFO
into its bounce buffer, \code{copy_to_user} fails on the unmapped page, and \code{read} returns
$-1$ with \code{errno} set to \code{EFAULT}, ``Bad address''. Note what the specification in
\secref{c5:app:b2} does with the bytes it had already taken: they are gone, consumed and never
delivered. (With the FIFO empty the driver returns \code{-EAGAIN} before it touches the pointer.)

\xpart{g} \code{cat} takes the 0 as end of file and exits with status 0, having printed nothing.
The loop ends at once, exactly as at the end of a file, and the program carries on as if the device
had finished. The \code{perror} check never fires, because 0 is not less than 0. \textbf{None of the
three reports a problem}, which is why returning 0 for ``nothing yet'' is a silent bug.

% ------------------------------------------------------------------------------------------
\solution[fillaccent3]{c5:ex:4}{A sequence of reads}{Hand calculation}

\xpart{a}

\begin{narrowtable}{@{}llrr@{}}
  \thead{Step} & \thead{Call} & \thead{Returns} & \thead{Level after} \\
  \midrule
  1 & \code{write(fd, "abcdefghij", 10)} & 8 & 8 \\
  2 & \code{read(fd, buf, 3)} & 3 & 5 \\
  3 & \code{write(fd, "klm", 3)} & 3 & 8 \\
  4 & \code{read(fd, buf, 100)} & 8 & 0 \\
  5 & \code{read(fd, buf, 100)} & \code{-EAGAIN} & 0 \\
\end{narrowtable}

\noindent Step 1 is a short write: \code{abcdefgh} goes in and \code{ij} is refused. Step 5 is
\code{-EAGAIN} and not 0, because the FIFO is empty rather than finished.

\xpart{b} \code{defghklm}: the five bytes step 2 left behind, then the three step 3 added, oldest
first. Step 4 overwrites the first eight bytes of \code{buf}, so the \code{abc} step 2 put there is
gone.

\xpart{c} Room is decided by the level, not by how much has ever been written. Step 2 took three
bytes out, leaving three slots free, and step 3 asks for exactly three; in a ring those are the
slots the first bytes occupied, reused by a write that wraps around the end of the storage.

\xpart{d} A FIFO of capacity 8 that wastes a slot holds 7, and two steps change. \textbf{Step 1}
returns 7 instead of 8, because only seven bytes fit. \textbf{Step 4} then returns 7 instead of 8,
because there are only seven to take, and \code{buf} holds \code{defgklm}: the \code{h} was never
stored. Steps 2, 3 and 5 return what they did before, 3, 3 and \code{-EAGAIN}, although the levels
after steps 1 to 3 are each one lower (7, 4, 7). This is the implementation that
\code{qa_fifo_uses_its_whole_capacity} exists to reject.

% ------------------------------------------------------------------------------------------
\solution[fillaccent4]{c5:ex:5}{Where does this belong}{Design}

\xpart{a} \code{read()} on the device node: it is the device's data, a stream consumed in order,
which is exactly what the file interface is for.

\xpart{b} A read-only sysfs attribute: one value per file, readable by any tool. It already exists
without further work, since the \code{capacity} module parameter appears as
\file{/sys/module/qa_chardev/parameters/capacity}.

\xpart{c} A read-only sysfs attribute. A field engineer's tools will read it on production
systems, where debugfs is usually not mounted and often not built in, so it has to live somewhere
that is guaranteed to be there.

\xpart{d} debugfs. It is only useful while you are developing the driver, and debugfs is the one
place that is explicitly not an interface: you may change it or delete it whenever you like.

\xpart{e} An \code{ioctl} on the open device. It is an operation rather than a value, and it acts
on the stream that the caller has open; terminals do the same thing through an \code{ioctl}, which
is what \code{tcflush} is. A write-only sysfs attribute is defensible, but sysfs is for attributes,
and an action dressed as one is harder to discover and to document.

\xpart{f} A read-only sysfs attribute: a single fact about the device that never changes, and one
that tools and udev rules can match on.

\xpart{g} The sysfs attributes and the \code{ioctl} are ABI: the moment a program or a script uses
them, their names, formats, numbers and meaning are fixed, and sysfs attributes are documented as
such under \file{Documentation/ABI/}. debugfs is explicitly not. What follows for a product is that
everything in sysfs or behind an \code{ioctl} has to be designed as if it will never change,
because it will not be allowed to, and that nothing a shipped tool or a field engineer depends on
may live in debugfs, which the product's kernel may not even contain. (The byte stream behind
\code{read()} is an interface too, but its contract, bytes in order, is the file interface's own
rather than one you invented.)

% ------------------------------------------------------------------------------------------
\solution[fillaccent4]{c5:ex:8}{Reading the tests for what they miss}{Design}

\xpart{a} Take the design \secref{c5:app:b1} suggests, a head and a level, with \code{qa_fifo_put}
copying to \code{(head + level) \% capacity} and then storing \code{level = level + n}. The FIFO has
capacity 8 and is empty, and two writers call \code{qa_fifo_put} with four bytes each, on two CPUs:
\begin{enumerate}
  \item Writer A reads \code{level}, 0, and copies \code{AAAA} into slots 0 to 3.
  \item Writer B reads \code{level}, still 0, and copies \code{BBBB} into the same slots,
    overwriting A's bytes.
  \item A stores \code{level = 0 + 4}; B stores \code{level = 0 + 4}.
\end{enumerate}
Both calls return 4, so both writers believe their bytes are in, but the FIFO holds four bytes,
all B's; A's data is lost without a trace. Change the timing and the level can end at 8 with slots
4 to 7 never written, so a reader is handed four bytes of garbage. \textbf{None of the ten tests
notices}, because every one of them runs on a single thread and never calls \code{qa_fifo_put}
twice at once.

\xpart{b} Any three of these, each worth testing:
\begin{itemize}
  \item \textbf{The \code{-EFAULT} path.} Nothing passes a bad pointer. Cheap to test: a program
    that calls \code{read} and \code{write} with an address in an unmapped page, expecting
    \code{EFAULT} and no oops.
  \item \textbf{Unloading while the device is open.} Nothing checks \code{.owner}. Cheap: hold the
    node open in one process and expect \code{rmmod} to refuse with ``Module is in use''.
  \item \textbf{A failing init.} Nothing makes any init step fail, so no error path is ever run.
    Test with \code{insmod qa_chardev.ko capacity=0}, which must fail cleanly because
    \code{qa_fifo_create(0)} returns \code{NULL}, leaving no node, no class and no number behind;
    later steps need fault injection.
  \item \textbf{Huge requests.} The tester writes 4096 bytes at a time. A \code{read} asking for a
    gigabyte must return a short count, not fail its bounce-buffer allocation.
  \item \textbf{Memory pressure.} The bounce buffer's allocation failing must give \code{-ENOMEM}
    rather than a crash; this needs the kernel's fault injection (\code{CONFIG_FAILSLAB}).
\end{itemize}

\xpart{c} Fewer than it seems, because a shell sees a command's exit status and its error message,
and those carry more than nothing: \code{cat} exits with status 0 on end of file and with 1 and
``Resource temporarily unavailable'' on \code{-EAGAIN}. Two checks it genuinely cannot make:
\begin{itemize}
  \item \textbf{\code{lseek} refused with \code{-ESPIPE}.} No command exposes the result of an
    \code{lseek}; \code{dd skip=} on a file it cannot seek quietly reads and discards instead.
  \item \textbf{Which error a call failed with.} The \code{-EAGAIN} and \code{-ENOSPC} checks
    compare \code{errno} with a constant. A script sees only that a tool failed and a message whose
    wording belongs to the tool and the C library, BusyBox's differing from coreutils', from tools
    that may retry a short write or treat one error like another before you ever see it.
\end{itemize}

\xpart{d} Neither. The KUnit suite never loads \code{qa_chardev}, and the tester and \file{run.sh}
only ever run the path on which every init step succeeds. What would catch it is a test that makes
\code{class_create} fail, for example by registering a class called \code{qa_fifo} first so that
it returns \code{-EEXIST}, run on a kernel with \code{CONFIG_DEBUG_KMEMLEAK}: after the failed
load, a kmemleak scan reports the FIFO's allocation as unreferenced.

% ------------------------------------------------------------------------------------------
\solution[fillaccent]{c5:ex:9}{A byte count worked out twice}{Cross-check}

\xpart{a} With a capacity of 256:

\begin{narrowtable}{@{}llrr@{}}
  \thead{Step} & \thead{Call} & \thead{Returns} & \thead{Level after} \\
  \midrule
  1 & write of 200 bytes & 200 & 200 \\
  2 & write of 100 bytes & 56 & 256 \\
  3 & read of 512 bytes & 256 & 0 \\
  4 & read of 512 bytes & \code{-EAGAIN} & 0 \\
  5 & write of 300 bytes & 256 & 256 \\
\end{narrowtable}

\noindent The device accepts $200 + 56 + 256 = 512$ bytes in all and delivers 256 of them, so 256
have passed all the way through and 256 are still inside when the sequence ends. A total of 600
counted the requests rather than the returns, and 500 counted step 2 as a failure.

\xpart{c} A driver that meets the specification gives exactly the table above, so any difference
is in your prediction or your driver. \textbf{Step 2} is the one people get wrong by predicting
\code{-ENOSPC}. The contract for \code{write} in \secref{c5:app:b2} decides it: \code{-ENOSPC} only
when the FIFO is \emph{full}, and otherwise as many bytes as fit, which here is $256 - 200 = 56$.
Predicting \code{-ENOSPC} there classes a value as an error; predicting 0 for step 4 classes an
error as end of file. The total counts step 2 as 56, as the machine does.

\xpart{d} With \code{capacity=128} the program sees 128, \code{-ENOSPC}, 128, \code{-EAGAIN} and
128. Steps 1, 3 and 5 each return 128 instead of 200, 256 and 256, and step 2 changes kind as well
as value: after step 1 the FIFO is full, so it is refused rather than short. Only step 4 is
unchanged, because it depends only on the FIFO being empty, and the 512-byte read in step 3 empties
a FIFO of any capacity up to 512.

\xpart{e} The shipped tester finds the capacity by writing until it is refused, so on the
\code{capacity=128} device it reports ``capacity measured from userspace: 128 bytes'' and passes
every check. Your program compares against 200, 56, 256, \code{-EAGAIN} and 256, and reports four
failures against a driver that is correct. The lesson: a test that encodes a configuration tests
the configuration. Derive the expectation from the system under test, by measuring it or by reading
the parameter back from sysfs, or a legitimate change of configuration turns into failures, and a
team that learns to shrug at failing tests stops seeing the real ones.

\xpart{f} Nothing is broken. The FIFO already held 200 of its 256 bytes, as it does after step 1,
so 56 bytes were free: the driver accepted 56 and said so. That is a short write, and their program
should have looked at the return value and written the remaining 244 bytes later, once a reader had
made room, in a loop, exactly as for a pipe or a socket. The line they had not read is the first row
of the table in \secref{c5:app:a:read-and-write}: this many bytes were transferred, and \textbf{a
short count is normal, not an error}.
